\chapter{Nếu Em Sai Thì Còn Ai Ở Lại Với Em}
\markboth{Nếu Em Sai Thì Còn Ai Ở Lại Với Em}{If I Am Wrong, Who Will Stay With Me}

\section{Sai Một Lần Là Bị Gắn Cả Năm}
\begin{truyen}{Sai Một Lần Là Bị Gắn Cả Năm}{One Mistake Becomes a Whole-Year Label}
\chuhoa{C}{ó những học sinh chỉ cần một lần quay cóp, một lần cãi, một lần hỗn là bị gắn luôn hình ảnh đó cho rất lâu.}
\textit{(Some students only need one act of cheating, arguing, or rudeness to carry that image for a long time.)}

Một lỗi thật sự đáng sửa.
\textit{(The mistake may truly need correction.)}

Nhưng nếu người lớn giữ nguyên cái nhìn đó quá lâu, đứa trẻ sẽ thấy mình không còn cơ hội làm mới nữa.
\textit{(But if adults keep that view too long, the child feels there is no chance to become new.)}

Có em nói rất thật: “Em sai một lần mà mọi người nhìn em như em sai cả đời.”
\textit{(Some students say very honestly, “I was wrong once, and everyone looks at me as if I'm wrong forever.”)}

Không có hy vọng sửa mình thì việc sửa mình cũng yếu đi.
\textit{(Without hope of renewal, the effort to change weakens.)}

\end{truyen}

\begin{baihoc}
Sửa lỗi là cần, nhưng để lại cửa quay về cũng cần không kém.
\textit{(Correction matters, but leaving a door back in matters just as much.)}
\end{baihoc}

\begin{ghinhoanh}
\textbf{Short English to remember:}

Correction matters, but so does leaving a way back.

\end{ghinhoanh}

\begin{tuhocnhanh}
\textbf{Gợi ý để dạy và tự nghĩ / Teaching and reflection prompts}

1. Vì sao nhãn cũ làm học sinh khó đổi? \textit{Why do old labels make change harder?}

2. Cửa quay về có nghĩa là gì trong giáo dục? \textit{What does a way back mean in education?}

3. Mình có đang giữ học sinh trong lỗi cũ quá lâu không? \textit{Am I keeping students trapped in old mistakes too long?}
\end{tuhocnhanh}
\ngancach

\section{Lúc Em Bị Cả Lớp Ghét}
\begin{truyen}{Lúc Em Bị Cả Lớp Ghét}{When the Whole Class Turns Against Me}
\chuhoa{C}{ó em làm sai thật và bị bạn bè đồng loạt quay lưng.}
\textit{(Some students truly do wrong and the whole class turns away from them.)}

Lúc đó, người lớn thường bận lo trật tự chung.
\textit{(At that moment, adults are often busy protecting the general order.)}

Nhưng ngay cả đứa làm sai cũng có thể đang chìm rất nhanh trong xấu hổ và hoảng.
\textit{(Yet even the student at fault may be sinking quickly into shame and panic.)}

Nếu không ai giúp em đi hết đoạn sửa sai đó, em rất dễ trượt sang lì hoặc bỏ mặc luôn chính mình.
\textit{(If no one helps them walk through the repair, they may become hardened or stop caring entirely.)}

\end{truyen}

\begin{baihoc}
Ngay cả khi bảo vệ tập thể, người lớn vẫn cần nhớ cách cứu một cá nhân đang bị đẩy ra rìa.
\textit{(Even while protecting the group, adults must remember how to rescue an individual pushed to the edge.)}
\end{baihoc}

\begin{ghinhoanh}
\textbf{Short English to remember:}

Even the student at fault still needs a path back.

\end{ghinhoanh}

\begin{tuhocnhanh}
\textbf{Gợi ý để dạy và tự nghĩ / Teaching and reflection prompts}

1. Vì sao học sinh làm sai vẫn cần được giữ? \textit{Why does a student at fault still need support?}

2. Điều gì dễ xảy ra nếu em bị đẩy ra hoàn toàn? \textit{What happens if the student is fully pushed out?}

3. Mình có thể giữ kỷ luật và giữ người cùng lúc thế nào? \textit{How can we protect discipline and the person at the same time?}
\end{tuhocnhanh}
\ngancach

\section{Xin Lỗi Xong Có Được Làm Lại Không}
\begin{truyen}{Xin Lỗi Xong Có Được Làm Lại Không}{After Apologizing, Am I Allowed to Begin Again}
\chuhoa{C}{ó những lời xin lỗi được chấp nhận bằng miệng nhưng không được chấp nhận bằng cách đối xử sau đó.}
\textit{(Some apologies are accepted in words but not in later treatment.)}

Người lớn nói đã tha rồi, nhưng cách nhìn vẫn cũ.
\textit{(Adults say they have forgiven, but their way of seeing remains unchanged.)}

Học sinh cảm được chuyện này rất nhanh.
\textit{(Students sense this very quickly.)}

Nếu xin lỗi mà không mở ra được một bắt đầu mới, lời xin lỗi sẽ dần mất giá trị trong mắt trẻ.
\textit{(If an apology does not open the door to a new start, it slowly loses meaning in the eyes of children.)}

\end{truyen}

\begin{baihoc}
Tha thứ thật sự luôn thể hiện một phần qua cách ta đối xử sau đó.
\textit{(Real forgiveness always shows itself partly in the way we treat someone afterward.)}
\end{baihoc}

\begin{ghinhoanh}
\textbf{Short English to remember:}

Real forgiveness appears in later treatment, not only in words.

\end{ghinhoanh}

\begin{tuhocnhanh}
\textbf{Gợi ý để dạy và tự nghĩ / Teaching and reflection prompts}

1. Vì sao trẻ cảm được sự tha thứ giả? \textit{Why can children sense fake forgiveness?}

2. Xin lỗi cần mở ra điều gì tiếp theo? \textit{What should an apology open up afterward?}

3. Mình có đang tha bằng lời nhưng giữ lỗi bằng mắt không? \textit{Am I forgiving with words while holding the fault in my eyes?}
\end{tuhocnhanh}
\ngancach

\section{Nếu Cả Người Lớn Cũng Mệt Với Em}
\begin{truyen}{Nếu Cả Người Lớn Cũng Mệt Với Em}{What If Even Adults Are Tired of Me}
\chuhoa{C}{ó những học sinh khó, và đúng là người lớn sẽ mệt thật.}
\textit{(Some students are difficult, and adults do genuinely get tired.)}

Nhưng điều học sinh rất sợ là nhìn thấy sự mệt đó chuyển thành thái độ muốn bỏ em cho xong.
\textit{(What students fear is seeing that fatigue turn into the attitude of wanting to give up on them.)}

Có em quậy hơn nữa để thử: liệu có ai còn giữ mình không.
\textit{(Some misbehave even more to test whether anyone will still hold them.)}

Những bài kiểm tra đau nhất của trẻ nhiều khi không nằm trong đề thi, mà nằm trong việc xem người lớn có còn ở lại không.
\textit{(A child's harshest tests often are not in exams, but in seeing whether adults will stay.)}

\end{truyen}

\begin{baihoc}
Nghề dạy không bắt người lớn không mệt. Nhưng nó đòi người lớn mệt mà vẫn cố không bỏ người.
\textit{(Teaching does not demand that adults never get tired. It asks that even in fatigue, they try not to abandon the child.)}
\end{baihoc}

\begin{ghinhoanh}
\textbf{Short English to remember:}

Teaching asks tired adults to keep trying not to give up on the child.

\end{ghinhoanh}

\begin{tuhocnhanh}
\textbf{Gợi ý để dạy và tự nghĩ / Teaching and reflection prompts}

1. Vì sao học sinh hay thử người lớn đúng lúc bị mệt nhất? \textit{Why do students often test adults when adults are most tired?}

2. Điều trẻ đang dò tìm là gì? \textit{What is the child trying to find out?}

3. Mình làm gì để vừa giữ mình vừa không bỏ em? \textit{What can I do to preserve myself without giving up on the student?}
\end{tuhocnhanh}
\ngancach

\section{Em Sai Rồi, Nhưng Đừng Nhìn Em Như Xong Rồi}
\begin{truyen}{Em Sai Rồi, Nhưng Đừng Nhìn Em Như Xong Rồi}{I Was Wrong, But Please Do Not Look at Me Like I Am Finished}
\chuhoa{C}{ó những ánh mắt của người lớn làm học sinh hiểu rất nhanh rằng chuyện của mình coi như hết cứu.}
\textit{(Some looks from adults tell students very quickly that their case is considered beyond saving.)}

Tôi từng nghe một em nói: “Em biết thầy cô giận. Em chỉ sợ thầy cô coi em xong rồi.”
\textit{(I once heard a student say, “I know teachers are angry. I am just afraid they think I am finished.”)}

Sai thì phải sửa.
\textit{(Wrongdoing must be corrected.)}

Nhưng nếu người lớn nhìn như thể không còn mùa nào sau lỗi này nữa, học sinh sẽ không còn muốn đi tiếp.
\textit{(But if adults look as if no season follows this mistake, students stop wanting to continue.)}

\end{truyen}

\begin{baihoc}
Sửa lỗi mà vẫn chừa lại hy vọng là cách sửa giúp trẻ còn muốn làm lại.
\textit{(Correcting while leaving hope intact helps children still want to begin again.)}
\end{baihoc}

\begin{ghinhoanh}
\textbf{Short English to remember:}

Correction works better when hope is left intact.

\end{ghinhoanh}

\begin{tuhocnhanh}
\textbf{Gợi ý để dạy và tự nghĩ / Teaching and reflection prompts}

1. Vì sao ánh mắt của người lớn lại nặng với học sinh như vậy? \textit{Why do adult looks carry so much weight?}

2. Hy vọng còn lại có vai trò gì sau khi trẻ sai? \textit{What role does hope play after a child has done wrong?}

3. Mình có đang sửa mà vô tình đóng luôn cánh cửa trở lại không? \textit{Am I correcting while accidentally shutting the way back?}
\end{tuhocnhanh}
\ngancach

\section{Em Bị Gọi Là Hư Nhiều Quá Nên Em Làm Theo Luôn}
\begin{truyen}{Em Bị Gọi Là Hư Nhiều Quá Nên Em Làm Theo Luôn}{I Was Called Bad So Often That I Started Acting Like It}
\chuhoa{C}{ó những học sinh bị gọi là hư, lì, bất trị suốt một thời gian dài.}
\textit{(Some students are called bad, stubborn, or hopeless for a very long time.)}

Ban đầu các em còn chống lại cái nhãn đó.
\textit{(At first they resist the label.)}

Sau nhiều lần, có em buông luôn: ai cũng nghĩ vậy thì thôi em sống vậy.
\textit{(After many repetitions, some give up: if everyone thinks so, I might as well become it.)}

Nhãn dán lâu ngày rất dễ thành vai diễn thật.
\textit{(A label repeated long enough easily becomes a real role.)}

\end{truyen}

\begin{baihoc}
Người lớn gọi tên học sinh thế nào có thể kéo các em gần hơn với điều đó.
\textit{(How adults name students can pull them closer to becoming that thing.)}
\end{baihoc}

\begin{ghinhoanh}
\textbf{Short English to remember:}

Repeated labels can become real roles.

\end{ghinhoanh}

\begin{tuhocnhanh}
\textbf{Gợi ý để dạy và tự nghĩ / Teaching and reflection prompts}

1. Vì sao nhãn lặp lại có thể thành bản sắc xấu? \textit{Why can repeated labels become identity?}

2. Học sinh buông xuôi ở điểm nào? \textit{At what point do students give up resisting?}

3. Mình đang gọi tên hành vi hay đang đóng khung cả con người? \textit{Am I naming a behavior or boxing in a whole person?}
\end{tuhocnhanh}
\ngancach

\section{Nếu Em Muốn Sửa Thì Có Ai Chờ Em Không}
\begin{truyen}{Nếu Em Muốn Sửa Thì Có Ai Chờ Em Không}{If I Want to Change, Will Anyone Wait for Me}
\chuhoa{C}{ó học sinh sửa không nhanh. Hôm nay đỡ, mai lại tái.}
\textit{(Some students do not change quickly. Better today, then slipping again tomorrow.)}

Người lớn rất dễ nản ở giai đoạn này.
\textit{(Adults easily get discouraged at this stage.)}

Nhưng sửa mình thật đâu có đi thành một đường thẳng đẹp.
\textit{(Real change rarely moves in a clean straight line.)}

Nhiều đứa trẻ chỉ cần biết vẫn có người chờ mình qua đoạn lặp lại ấy.
\textit{(Many children only need to know that someone will stay through the repeated falling.)}

\end{truyen}

\begin{baihoc}
Kiên nhẫn với quá trình sửa của học sinh là phần khó nhưng rất thật của giáo dục.
\textit{(Patience with a student's uneven process of change is a difficult but very real part of education.)}
\end{baihoc}

\begin{ghinhoanh}
\textbf{Short English to remember:}

Real change is rarely a straight line.

\end{ghinhoanh}

\begin{tuhocnhanh}
\textbf{Gợi ý để dạy và tự nghĩ / Teaching and reflection prompts}

1. Vì sao người lớn hay nản khi học sinh tái phạm? \textit{Why do adults get discouraged when students relapse?}

2. Một quá trình sửa thật thường trông như thế nào? \textit{What does a real process of change often look like?}

3. Mình có đang đòi học sinh tiến bộ quá thẳng không? \textit{Am I expecting progress to be too linear?}
\end{tuhocnhanh}
\ngancach

\section{Em Xin Lỗi Rồi, Nhưng Em Vẫn Thấy Mọi Người Né Em}
\begin{truyen}{Em Xin Lỗi Rồi, Nhưng Em Vẫn Thấy Mọi Người Né Em}{I Apologized, But People Still Avoid Me}
\chuhoa{C}{ó em đã xin lỗi thật, đã nhận sai thật, nhưng đi trong lớp vẫn thấy không khí khác hẳn.}
\textit{(Some students truly apologize and admit fault, yet still feel a changed atmosphere around them in class.)}

Các bạn né. Người lớn dè.
\textit{(Friends avoid. Adults stay cautious.)}

Lúc đó học sinh sẽ tự hỏi: vậy xin lỗi để làm gì.
\textit{(Then students begin to ask: what was the apology even for.)}

Tha thứ không có nghĩa là quên ngay, nhưng nếu không có dấu hiệu nào của chỗ trở lại, trẻ sẽ mất niềm tin vào việc sửa sai.
\textit{(Forgiveness does not mean instant forgetting, but without any sign of return, children lose faith in repair.)}

\end{truyen}

\begin{baihoc}
Sau lời xin lỗi, lớp học cũng cần học cách cho người sai một con đường quay lại.
\textit{(After an apology, the classroom must also learn how to offer a path back.)}
\end{baihoc}

\begin{ghinhoanh}
\textbf{Short English to remember:}

Repair needs a path back, not just an apology.

\end{ghinhoanh}

\begin{tuhocnhanh}
\textbf{Gợi ý để dạy và tự nghĩ / Teaching and reflection prompts}

1. Vì sao học sinh dễ tuyệt vọng sau khi xin lỗi mà vẫn bị né? \textit{Why do students lose hope when they apologize and are still avoided?}

2. Cả lớp cần học gì sau một lỗi lớn? \textit{What does the whole class need to learn after a major mistake?}

3. Mình có đang dạy đủ về quay lại sau đổ vỡ không? \textit{Am I teaching enough about return after failure?}
\end{tuhocnhanh}
\ngancach

\section{Người Lớn Mệt Với Em Không Sao, Chỉ Xin Đừng Bỏ Em Trong Sự Mệt Đó}
\begin{truyen}{Người Lớn Mệt Với Em Không Sao, Chỉ Xin Đừng Bỏ Em Trong Sự Mệt Đó}{Adults May Be Tired of Me, but Please Do Not Leave Me There}
\chuhoa{C}{ó lần một học sinh nói đúng một câu: “Em biết thầy cô mệt vì em. Nhưng đừng bỏ em nha.”}
\textit{(Once a student said only this: “I know teachers are tired because of me. But please don't give up on me.”)}

Nghe xong không thể coi sự quậy chỉ là sự quậy nữa.
\textit{(After hearing that, it becomes impossible to see the misbehavior as mere misbehavior.)}

Đằng sau nhiều đứa trẻ khó là một nỗi sợ bị bỏ rất thật.
\textit{(Behind many difficult children is a very real fear of abandonment.)}

Người lớn không cần giả vờ mình không mệt. Nhưng cần biết mệt mà vẫn giữ được người.
\textit{(Adults do not need to pretend they are never tired. But they need to know how to stay while tired.)}

\end{truyen}

\begin{baihoc}
Điều khó trong giáo dục không phải là không mệt, mà là không để sự mệt biến thành bỏ mặc.
\textit{(The difficulty in education is not avoiding fatigue, but refusing to let fatigue become abandonment.)}
\end{baihoc}

\begin{ghinhoanh}
\textbf{Short English to remember:}

The real challenge is not fatigue, but refusing to let fatigue become abandonment.

\end{ghinhoanh}

\begin{tuhocnhanh}
\textbf{Gợi ý để dạy và tự nghĩ / Teaching and reflection prompts}

1. Vì sao câu xin đừng bỏ em lại nặng như vậy? \textit{Why is the plea “don't give up on me” so heavy?}

2. Mệt và bỏ mặc khác nhau ở đâu? \textit{How is fatigue different from abandonment?}

3. Mình giữ học sinh bằng cách nào trong lúc bản thân cũng mệt? \textit{How do I keep holding the student while tired myself?}
\end{tuhocnhanh}
\ngancach

\section{Có Khi Điều Giữ Em Lại Chỉ Là Một Người Chưa Chốt Sổ Với Em}
\begin{truyen}{Có Khi Điều Giữ Em Lại Chỉ Là Một Người Chưa Chốt Sổ Với Em}{Sometimes What Saves Me Is One Adult Who Has Not Closed the Book on Me}
\chuhoa{C}{ó những học sinh còn đi học tiếp không phải vì mọi thứ đã ổn, mà vì vẫn còn một người chưa hết tin.}
\textit{(Some students keep coming to school not because everything is okay, but because one person has not stopped believing.)}

Một giáo viên chủ nhiệm. Một cô giám thị. Một người bảo vệ hay chào em.
\textit{(A homeroom teacher. A supervisor. A guard who still greets the child.)}

Nhiều khi đời học trò được giữ lại bởi đúng một mắt xích như vậy.
\textit{(Sometimes a school life is held together by exactly one such link.)}

Giữa lúc cả thế giới của em như đang chốt sổ, chỉ cần còn một người chưa chốt là em còn đường trở lại.
\textit{(When the whole world seems ready to close the file, one person who does not can leave a path back.)}

\end{truyen}

\begin{baihoc}
Trong đời một học sinh lạc, một người chưa bỏ cuộc có thể có giá trị lớn hơn nhiều bài giảng đúng.
\textit{(In the life of a drifting student, one adult who has not given up can matter more than many correct lectures.)}
\end{baihoc}

\begin{ghinhoanh}
\textbf{Short English to remember:}

One adult who has not given up can keep a student in school.

\end{ghinhoanh}

\begin{tuhocnhanh}
\textbf{Gợi ý để dạy và tự nghĩ / Teaching and reflection prompts}

1. Vì sao chỉ một người còn tin cũng có thể đủ để giữ học sinh lại? \textit{Why can one believing adult be enough to keep a student going?}

2. Vai trò của người đó khác gì với các biện pháp chung? \textit{How is that role different from general school measures?}

3. Mình có từng là người chưa chốt sổ với một học sinh nào không? \textit{Have I ever been the adult who did not close the book on a student?}
\end{tuhocnhanh}
